% Options for packages loaded elsewhere
\PassOptionsToPackage{unicode}{hyperref}
\PassOptionsToPackage{hyphens}{url}
\documentclass[
  12pt,
]{ctexart}
\usepackage{xcolor}
\usepackage{amsmath,amssymb}
\setcounter{secnumdepth}{-\maxdimen} % remove section numbering
\usepackage{iftex}
\ifPDFTeX
  \usepackage[T1]{fontenc}
  \usepackage[utf8]{inputenc}
  \usepackage{textcomp} % provide euro and other symbols
\else % if luatex or xetex
  \usepackage{unicode-math} % this also loads fontspec
  \defaultfontfeatures{Scale=MatchLowercase}
  \defaultfontfeatures[\rmfamily]{Ligatures=TeX,Scale=1}
\fi
\usepackage{lmodern}
\ifPDFTeX\else
  % xetex/luatex font selection
\fi
% Use upquote if available, for straight quotes in verbatim environments
\IfFileExists{upquote.sty}{\usepackage{upquote}}{}
\IfFileExists{microtype.sty}{% use microtype if available
  \usepackage[]{microtype}
  \UseMicrotypeSet[protrusion]{basicmath} % disable protrusion for tt fonts
}{}
\makeatletter
\@ifundefined{KOMAClassName}{% if non-KOMA class
  \IfFileExists{parskip.sty}{%
    \usepackage{parskip}
  }{% else
    \setlength{\parindent}{0pt}
    \setlength{\parskip}{6pt plus 2pt minus 1pt}}
}{% if KOMA class
  \KOMAoptions{parskip=half}}
\makeatother
\usepackage{longtable,booktabs,array}
\usepackage{calc} % for calculating minipage widths
% Correct order of tables after \paragraph or \subparagraph
\usepackage{etoolbox}
\makeatletter
\patchcmd\longtable{\par}{\if@noskipsec\mbox{}\fi\par}{}{}
\makeatother
% Allow footnotes in longtable head/foot
\IfFileExists{footnotehyper.sty}{\usepackage{footnotehyper}}{\usepackage{footnote}}
\makesavenoteenv{longtable}
\setlength{\emergencystretch}{3em} % prevent overfull lines
\providecommand{\tightlist}{%
  \setlength{\itemsep}{0pt}\setlength{\parskip}{0pt}}
\usepackage{bookmark}
\IfFileExists{xurl.sty}{\usepackage{xurl}}{} % add URL line breaks if available
\urlstyle{same}
\hypersetup{
  hidelinks,
  pdfcreator={LaTeX via pandoc}}

\author{SCX}
\date{2026}

\begin{document}

\section{竞争分析报告：数据估值与主动学习}\label{ux7adeux4e89ux5206ux6790ux62a5ux544aux6570ux636eux4f30ux503cux4e0eux4e3bux52a8ux5b66ux4e60}

\begin{quote}
SCX (State-Conditioned eXpertise) 框架的竞争格局全景分析
覆盖时间范围：2023-2026 生成日期：2026-06-25
\end{quote}

\begin{center}\rule{0.5\linewidth}{0.5pt}\end{center}

\subsection{目录}\label{ux76eeux5f55}

\begin{enumerate}
\def\labelenumi{\arabic{enumi}.}
\tightlist
\item
  \hyperref[1-ux65b9ux6cd5ux8bbaux8bf4ux660e]{方法论说明}
\item
  \hyperref[a-data-valuation-ux65b9ux6cd5]{A. Data Valuation 方法}
\item
  \hyperref[b-llm-ux6570ux636eux9009ux62e9]{B. LLM 数据选择}
\item
  \hyperref[c-active-learning]{C. Active Learning}
\item
  \hyperref[d-ux7adeux4e89ux683cux5c40ux603bux7ed3]{D. 竞争格局总结}
\item
  \hyperref[e-scx-ux7684ux72ecux7279ux5b9aux4f4d]{E. SCX 的独特定位}
\item
  \hyperref[7-ux53c2ux8003ux6587ux732e]{参考文献}
\end{enumerate}

\begin{center}\rule{0.5\linewidth}{0.5pt}\end{center}

\subsection{1. 方法论说明}\label{ux65b9ux6cd5ux8bbaux8bf4ux660e}

本报告通过系统化 Web 搜索 + 论文检索 + 源码分析三种渠道，对四大领域的近
30
种方法进行逐一评估。每个方法按照统一模板分析：解决什么问题、数学对象是什么、与
SCX 的核心差异、SCX 的相对优势。

SCX 框架的核心定义（供对手比较使用）：

\begin{itemize}
\tightlist
\item
  \textbf{专家可靠性}：\texttt{SCX\_m(s)\ =\ P(ℓ(f\_m(x),\ y)\ \textless{}\ τ\ \textbar{}\ x\ ∈\ s)}
  --- 不是全局常数，而是\textbf{状态条件函数}
\item
  \textbf{数据价值}：\texttt{V(s)\ =\ r̄(s)\ ·\ ρ(s)\ ·\ L(s)\ ·\ {[}1\ -\ D(s){]}\ ·\ max\_m\ SCX\_m(s)}
  --- 不是样本固有属性
\item
  \textbf{状态划分}：\texttt{s:\ X\ →\ S}，状态空间由学习/聚类自动发现
\item
  \textbf{两类获取}：state-wise acquisition（先选高价值状态）和
  point-wise acquisition（状态内选代表点）
\end{itemize}

\begin{center}\rule{0.5\linewidth}{0.5pt}\end{center}

\subsection{A. Data Valuation
方法}\label{a.-data-valuation-ux65b9ux6cd5}

\subsubsection{A.1 Data Shapley
及其变体}\label{a.1-data-shapley-ux53caux5176ux53d8ux4f53}

\paragraph{A.1.1 经典 Data
Shapley}\label{a.1.1-ux7ecfux5178-data-shapley}

\begin{itemize}
\tightlist
\item
  \textbf{论文}：Ghorbani \& Zou, ``Data Shapley: Equitable Valuation of
  Data for Machine Learning'', ICML 2019
\item
  \textbf{链接}：https://arxiv.org/abs/1904.02868
\end{itemize}

\textbf{解决问题}：对每个训练样本赋予一个公平的''价值分数''，衡量其对模型性能的边际贡献。

\textbf{数学对象}：合作博弈论中的 Shapley 值。对样本 i：

\begin{verbatim}
φ_i = (1/N!) Σ_{π∈Π} [v(S_π(i) ∪ {i}) - v(S_π(i))]
\end{verbatim}

其中 v(S) 是训练集 S 上训练模型的性能，Π 是所有排列，S\_π(i) 是排列 π 中
i 之前的样本集。

\textbf{与 SCX 的核心差异}： 1. \textbf{全局价值 vs 条件价值}：Data
Shapley 给每个样本一个全局标量；SCX
的价值是状态条件的------一个样本在状态 s
中的价值取决于当前模型在该状态的缺陷和专家的可靠性。 2. \textbf{单模型
vs 多专家}：Data Shapley 基于一个模型 v(S)；SCX 天然处理多个专家。 3.
\textbf{计算开销}：Data Shapley 精确计算 O(2\^{}N)，MC 近似仍需要 O(N²)
次重训练。

\textbf{SCX 相对优势}： - Data Shapley
无法区分''一个困难状态的样本''和''一个噪声样本''------两者都可以有高
Shapley 值。SCX 通过状态内部一致性 C(s) 做区分。 - Data Shapley
无法利用专家先验知识------它对每个样本的估值从头计算。SCX
可用专家条件可靠性加速估值。 - Data Shapley 的估值对模型重训练敏感，SCX
的状态划分更稳定。

\begin{center}\rule{0.5\linewidth}{0.5pt}\end{center}

\paragraph{A.1.2 截断 Monte Carlo Shapley
(TMC-Shapley)}\label{a.1.2-ux622aux65ad-monte-carlo-shapley-tmc-shapley}

\begin{itemize}
\tightlist
\item
  \textbf{论文}：Ghorbani \& Zou, ICML 2019（同上）
\item
  \textbf{变体}：对 Shapley 值的 MC
  估计进行截断：当边际贡献收敛时提前终止排列采样
\end{itemize}

\textbf{改进}：计算量从 O(N²) 降为约 O(NlogN)，仍需要大量重训练。

\textbf{SCX 差异}：TMC 只是计算加速策略，不改变全局价值的基本逻辑。SCX
的 state-wise valuation 天然可并行（各状态独立估值），且不需要排列。

\begin{center}\rule{0.5\linewidth}{0.5pt}\end{center}

\paragraph{A.1.3 KNN-Shapley}\label{a.1.3-knn-shapley}

\begin{itemize}
\tightlist
\item
  \textbf{论文}：Jia et al., ``Towards Efficient Data Valuation Based on
  the Shapley Value'', AISTATS 2019
\item
  \textbf{链接}：https://arxiv.org/abs/1902.10275
\end{itemize}

\textbf{数学对象}：用 KNN 作为代理模型（surrogate model），利用 KNN
的固有性质实现精确 O(n log n) 的 Shapley 值计算。

\textbf{改进}：首个实用大规模数据估值方法，不需要大量重训练。

\textbf{SCX 差异}： - KNN-Shapley
将估值绑定到一个特定代理模型（KNN）。SCX
不依赖特定模型，而是使用状态划分的密度和一致性。 - KNN
距离在高维空间失效（curse of dimensionality），SCX
的状态空间可以通过降维 + 聚类自动构造。

\begin{center}\rule{0.5\linewidth}{0.5pt}\end{center}

\paragraph{A.1.4 CHG Shapley}\label{a.1.4-chg-shapley}

\begin{itemize}
\tightlist
\item
  \textbf{论文}：Li et al., ``CHG Shapley: Efficient Data Valuation and
  Selection towards Trustworthy Machine Learning'', 2024
\item
  \textbf{链接}：https://arxiv.org/abs/2406.11730
\end{itemize}

\textbf{数学对象}：提出 Compound Hardness \& Gradient (CHG)
效用函数，在每个 epoch 近似子集效用，Shapley 值有闭式解。

\textbf{改进}：相比 marginal contribution
方法有二次加速。适用于标签噪声检测和类别不平衡。

\textbf{SCX 差异}：CHG 仍然输出全局标量值。SCX 的 state-conditioning
能提供更细粒度的分析------例如在状态 A 中某个样本是噪声，在状态 B
中同一样本可能是关键样本。

\begin{center}\rule{0.5\linewidth}{0.5pt}\end{center}

\paragraph{A.1.5 Fast-DataShapley}\label{a.1.5-fast-datashapley}

\begin{itemize}
\tightlist
\item
  \textbf{论文}：``Fast-DataShapley: Neural Modeling for Training Data
  Valuation'', ACM WSDM 2026
\item
  \textbf{链接}：https://arxiv.org/abs/2506.05281
\end{itemize}

\textbf{数学对象}：用加权最小二乘特征化 Shapley
值，训练一个一次可重用的解释器模型。不需要为新的测试样本重新估值。

\textbf{改进}：2 倍性能提升，训练加速两个数量级。

\textbf{SCX 差异}：Fast-DataShapley 加速了传统 Shapley
计算，但仍是\textbf{事后解释}------需要先训练好模型再估值。SCX
是\textbf{主动框架}，在训练/标注过程中决定哪些数据最有价值。

\begin{center}\rule{0.5\linewidth}{0.5pt}\end{center}

\paragraph{A.1.6 Asymmetric Data Shapley
(ADS)}\label{a.1.6-asymmetric-data-shapley-ads}

\begin{itemize}
\tightlist
\item
  \textbf{论文}：``Rethinking Data Value: Asymmetric Data Shapley for
  Structure-Aware Valuation'', 2025
\item
  \textbf{链接}：https://arxiv.org/abs/2511.12863
\end{itemize}

\textbf{数学对象}：放松经典 Shapley
的对称公理，引入方向性/时序依赖（数据增广、联邦学习、多阶段微调）。

\textbf{改进}：保持效率性、线性性和集群对称性。

\textbf{SCX 差异}：ADS 的''asymmetric''是指样本之间的关系不对称（如 A
增广出 B 则 A 价值应大于 B）。SCX
的''state-conditioned''是指价值随状态变化------两个正交的改进方向，可以结合。

\begin{center}\rule{0.5\linewidth}{0.5pt}\end{center}

\paragraph{A.1.7 Geometric Data Valuation via Leverage
Scores}\label{a.1.7-geometric-data-valuation-via-leverage-scores}

\begin{itemize}
\tightlist
\item
  \textbf{论文}：``Geometric Data Valuation via Leverage Scores'', 2025
\item
  \textbf{链接}：https://arxiv.org/abs/2511.02100
\end{itemize}

\textbf{数学对象}：非 Shapley 的几何替代方案，使用统计 leverage
score。满足 dummy、efficiency 和 symmetry 公理。

\textbf{与 SCX 联系}：Leverage scores
本质上衡量了样本在特征空间中的''孤立度''------这近似于 SCX 中密度 ρ(s)
的一个特例。但 SCX
的状态空间更丰富：不仅包含几何密度，还包含标签一致性、专家可靠性和可学习性。

\textbf{SCX
优势}：几何方法只用到特征空间的几何，忽略了标签信息、专家知识和噪声结构。

\begin{center}\rule{0.5\linewidth}{0.5pt}\end{center}

\subsubsection{A.2 Influence Function
方法}\label{a.2-influence-function-ux65b9ux6cd5}

\paragraph{A.2.1 经典 Influence
Functions}\label{a.2.1-ux7ecfux5178-influence-functions}

\begin{itemize}
\tightlist
\item
  \textbf{原著}：Cook \& Weisberg, ``Residuals and Influence in
  Regression'', 1982
\item
  \textbf{ML 引入}：Koh \& Liang, ``Understanding Black-box Predictions
  via Influence Functions'', ICML 2017
\item
  \textbf{链接}：https://arxiv.org/abs/1703.04730
\end{itemize}

\textbf{数学对象}：通过 Hessian 矩阵追溯模型预测对训练样本的依赖：

\begin{verbatim}
I(x_i, x_test) = -∇_θ L(x_test)^T H_θ^{-1} ∇_θ L(x_i)
\end{verbatim}

\textbf{解决问题}：回答''删除这个训练点会如何改变测试点的预测？''

\textbf{与 SCX 的核心差异}： 1. \textbf{单模型 vs 多专家}：Influence
Functions 分析单一模型的训练数据影响；SCX 分析多个专家在状态上的可靠性。
2. \textbf{局部 vs 全局}：IF 是逐样本-逐测试点的；SCX 是状态级的。 3.
\textbf{Hessian 计算}：IF 需要高效的 Hessian 逆近似（LiSSA,
CG），对大型模型极其昂贵。

\textbf{SCX 相对优势}： - IF 无法处理多个专家之间的路由决策问题。 - IF
的估值是测试点依赖的（需要知道 x\_test），SCX
的价值是状态依赖的（不需要特定的测试点）。 - IF
对模型假设敏感（需要损失函数二阶可微、模型收敛到局部极小）。

\begin{center}\rule{0.5\linewidth}{0.5pt}\end{center}

\paragraph{A.2.2 TRAK (Tracing with the Randomly-Projected After
Kernel)}\label{a.2.2-trak-tracing-with-the-randomly-projected-after-kernel}

\begin{itemize}
\tightlist
\item
  \textbf{论文}：Park et al., ``TRAK: Attributing Model Behavior at
  Scale'', ICML 2023
\item
  \textbf{链接}：https://arxiv.org/abs/2212.02522
\item
  \textbf{代码}：https://github.com/JD-ETH/trak
\end{itemize}

\textbf{数学对象}：随机投影的 NTK（Neural Tangent
Kernel）版本，将梯度投影到低维空间后计算影响：

\begin{verbatim}
τ(x_i, x) = ⟨ϕ(x_i), M^{-1} ϕ(x)⟩
\end{verbatim}

其中 ϕ(x\_i) = J\_θ(x\_i) · P，P 是随机投影矩阵，M 是经验 NTK。

\textbf{改进}：比影响函数便宜 2-3 个数量级，可扩展到大型视觉和语言模型。

\textbf{SCX 差异}：TRAK 仍然是事后归因方法，不是主动数据选择框架。TRAK
估值依赖于训练好的模型和投影矩阵。SCX
不依赖训练好的模型------它基于状态划分和专家先验。

\begin{center}\rule{0.5\linewidth}{0.5pt}\end{center}

\paragraph{A.2.3 D-TRAK}\label{a.2.3-d-trak}

\begin{itemize}
\tightlist
\item
  \textbf{论文}：Zheng et al., ``Intriguing Properties of Data
  Attribution on Diffusion Models'', ICLR 2024
\item
  \textbf{链接}：https://github.com/qwkkkkk/D-TRAK
\end{itemize}

\textbf{数学对象}：TRAK
在扩散模型上的适配，对理论上不合理的设计选择在实践中表现出色。

\textbf{SCX 差异}：领域特定（扩散模型），不涉及状态条件化或多专家。

\begin{center}\rule{0.5\linewidth}{0.5pt}\end{center}

\paragraph{A.2.4 LoGra / LogIX}\label{a.2.4-logra-logix}

\begin{itemize}
\tightlist
\item
  \textbf{论文}：Choe et al., ``What is Your Data Worth to GPT?
  LLM-Scale Data Valuation with Influence Functions'', NeurIPS 2025
\item
  \textbf{链接}：https://arxiv.org/abs/2405.13954
\end{itemize}

\textbf{数学对象}：利用梯度的 Kronecker 乘积结构（类似
LoRA），将梯度投影复杂度从 O(nk) 降到 O(√(nk))。LoGra 算法 + LogIX
软件框架。

\textbf{改进}：针对 Llama3-8B 在 1B token 数据上实现最高 6500x
吞吐量提升，5x GPU 内存降低。

\textbf{SCX 差异}：LoGra 专注于让传统影响函数在 LLM
规模上可行，不改变影响函数的基本范式（单模型、事后归因、逐样本估值）。SCX
的 state-conditioned 视角在方法论上正交。

\begin{center}\rule{0.5\linewidth}{0.5pt}\end{center}

\paragraph{A.2.5 Generalized Group Data Attribution
(GGDA)}\label{a.2.5-generalized-group-data-attribution-ggda}

\begin{itemize}
\tightlist
\item
  \textbf{论文}：``Generalized Group Data Attribution'', 2024
\item
  \textbf{链接}：https://arxiv.org/abs/2410.09940
\end{itemize}

\textbf{数学对象}：将 TRAK/IF/TracIn
从单个训练点扩展到\textbf{训练点组}。10x-50x
加速，适用于数据剪枝和噪声标签识别。

\textbf{SCX 差异}：GGDA
的''组''是任意预定义的（如属于同一类的所有样本），而 SCX
的''状态''是通过误差/行为模式\textbf{自动发现}的语义结构。GGDA
的分组先验是人工给定的，SCX 的状态可以从数据中学习。

\begin{center}\rule{0.5\linewidth}{0.5pt}\end{center}

\subsubsection{A.3 LOO / 留出法}\label{a.3-loo-ux7559ux51faux6cd5}

\paragraph{A.3.1 经典
Leave-One-Out}\label{a.3.1-ux7ecfux5178-leave-one-out}

\textbf{数学对象}：

\begin{verbatim}
LOO_i = v(D) - v(D_{-i})
\end{verbatim}

即训练全部数据 vs 去掉样本 i 后的性能差异。

\textbf{问题}： - O(N) 次重训练，实际不可行 - 在大数据集上 LOO
值普遍趋近于零（单个样本影响小） - 不满足公平性公理（副本值不为零）

\paragraph{A.3.2 Kairos}\label{a.3.2-kairos}

\begin{itemize}
\tightlist
\item
  \textbf{论文}：``Kairos: Scalable Model-Agnostic Data Valuation'',
  2025
\item
  \textbf{链接}：https://arxiv.org/abs/2506.23799
\end{itemize}

\textbf{数学对象}：基于 MMD 的 LOO 近似，不需要重训练，支持在线增量更新
O(mN)。误差 O(1/N²)，比 Data Shapley 快 50 倍。

\textbf{SCX 差异}：Kairos
仍是逐样本标量估值，不提供状态条件化或多专家路由能力。

\begin{center}\rule{0.5\linewidth}{0.5pt}\end{center}

\subsubsection{A.4 pyDVL 生态}\label{a.4-pydvl-ux751fux6001}

\begin{itemize}
\tightlist
\item
  \textbf{项目}：https://github.com/aai-institute/pyDVL
\item
  \textbf{文档}：https://pydvl.org
\item
  \textbf{版本}：v0.10.0 (2024)
\end{itemize}

\textbf{实现功能}： - Shapley 值（精确、MC
排列、截断、群测试、class-wise） - Banzhaf 指数（含 MSR Banzhaf） - Beta
Shapley 等半值 - Least Core（精确和 MC） - LOO、Data-OOB、KNN
Shapley、Owen 采样 - 影响函数（DirectInfluence、CG、Block
CG、LiSSA、Arnoldi、EKFAC、Nystrom、DataInf）

\textbf{与 SCX 的关系}：pyDVL
是目前最完整的开源数据估值工具库，但它是一个库------不是方法论。SCX
可以\textbf{使用 pyDVL
作为子模块}来计算某些组件（如状态内的价值排序），但 pyDVL
本身不提供状态划分、多专家路由或 state-wise acquisition。

\begin{center}\rule{0.5\linewidth}{0.5pt}\end{center}

\subsection{B. LLM 数据选择}\label{b.-llm-ux6570ux636eux9009ux62e9}

\subsubsection{B.1 DataComp-LM (DCLM)}\label{b.1-datacomp-lm-dclm}

\begin{itemize}
\tightlist
\item
  \textbf{论文}：``DataComp-LM: In search of the next generation of
  training sets for language models'', NeurIPS 2024
\item
  \textbf{链接}：https://arxiv.org/abs/2406.11794
\item
  \textbf{网站}：https://datacomp.ai/dclm/
\end{itemize}

\textbf{解决问题}：为语言模型训练提供标准化的数据策展
benchmark------固定训练代码和预算，让参与者比拼数据策展策略。

\textbf{核心发现}：基于模型的过滤（model-based
filtering）对高质量训练集构建至关重要。DCLM-Baseline (7B) 在 MMLU 上达到
64\%，用 2.6T token 匹配 Mistral-7B（63\%）和 Llama 3
8B（66\%），但计算量减少 6.6 倍。

\textbf{与 SCX 的关系}：DCLM 是 benchmark
平台，不是数据选择算法本身。但其核心发现------\textbf{数据质量的重要性超过模型架构}------恰恰支持了
SCX 的假设。SCX 可以作为一种''智能数据策展策略''参与 DCLM 的过滤
track。DCLM 的 53 个下游评估任务可作为 SCX
状态划分效果的大规模验证平台。

\textbf{SCX 差异与优势}： - DCLM 研究的是预训练数据选择（海量 web
文本过滤）；SCX 主要针对微调/下游任务（指令数据选择、主动学习采样） -
DCLM 的方法（perplexity 过滤、heuristic 过滤）是\textbf{无监督的}；SCX
利用\textbf{监督信号}（专家误差、标签一致性） - SCX 的状态划分可能比
DCLM
的简单过滤更精细化------不仅过滤掉''坏''数据，还识别''什么状态下需要什么类型的数据''

\begin{center}\rule{0.5\linewidth}{0.5pt}\end{center}

\subsubsection{B.2 LESS (Low-rank gradiEnt Similarity
Search)}\label{b.2-less-low-rank-gradient-similarity-search}

\begin{itemize}
\tightlist
\item
  \textbf{论文}：Xia et al., ``LESS: Selecting Influential Data for
  Targeted Instruction Tuning'', ICML 2024
\item
  \textbf{链接}：https://arxiv.org/abs/2402.04333
\item
  \textbf{代码}：https://github.com/princeton-nlp/LESS
\end{itemize}

\textbf{解决问题}：针对目标能力（如推理、翻译）从大规模指令数据中选出最相关的微调样本。

\textbf{数学对象}： 1. 在部分数据上训练小 LoRA adapter → 计算 Adam LoRA
梯度特征 → 构建梯度数据存储 2. 对目标任务（少量验证集）计算梯度特征 →
通过内积/余弦相似度选择最相似训练样本 3. 关键适配：使影响函数与 Adam
优化器和变长指令数据兼容

\textbf{关键结果}： - 5\% 的 LESS 筛选数据往往优于全部数据 -
小模型（LLaMA-2-7B）选的数据可迁移到大模型（LLaMA-2-13B, MISTRAL-7B） -
优于随机选择 2-5 个百分点（MMLU, TyDiQA, BBH）

\textbf{计算约束分析}（后续工作，arXiv:2410.16208）： -
固定训练预算：LESS 和 perplexity 方法最优 -
固定总计算预算（含选择成本）：BM25 和 embedding 方法在 7B-13B 规模更优 -
70B 规模：LESS 最优（选择成本被训练成本摊销）

\textbf{与 SCX 的核心差异}： 1. \textbf{梯度空间 vs 状态空间}：LESS
在梯度空间做相似性匹配；SCX
在语义/状态空间做划分。梯度空间维度高且解释性差，状态空间可解释且稳定。
2. \textbf{单验证集 vs 多状态}：LESS
的目标是从一个验证任务出发筛选数据；SCX
可以同时处理多个状态（即多个子任务/模式）。 3. \textbf{专家的缺失}：LESS
只用一个小模型计算梯度相似性；SCX 利用多个专家的（先验）可靠性。

\textbf{SCX 相对优势}： - LESS
不能区分''高误差是因为样本困难''还是''高误差是因为样本噪声''------这是
SCX 通过状态内一致性 C(s) 解决的问题。 - LESS 需要预训练 LoRA adapter
来计算梯度，SCX
的状态划分可以来自先验知识（如材料科学中的键类型）或聚类。 - LESS
的验证集假设（有少量目标任务数据）在某些场景不成立；SCX 的 state-value
可以完全无验证集（通过 r̄(s)·ρ(s)·L(s)·{[}1-D(s){]}
自动发现高价值状态）。

\begin{center}\rule{0.5\linewidth}{0.5pt}\end{center}

\subsubsection{B.3 S2L (SmallToLarge)}\label{b.3-s2l-smalltolarge}

\begin{itemize}
\tightlist
\item
  \textbf{论文}：Yang et al., ``SmallToLarge (S2L): Scalable Data
  Selection for Fine-tuning Large Language Models by Summarizing
  Training Trajectories of Small Models'', NeurIPS 2024
\item
  \textbf{链接}：https://arxiv.org/abs/2403.07384
\item
  \textbf{代码}：https://github.com/BigML-CS-UCLA/S2L
\end{itemize}

\textbf{解决问题}：用比目标模型小 40-100
倍的小模型训练轨迹来指导大模型的数据选择。

\textbf{数学对象}： 1. 在全部数据上训练小模型（如 Pythia-70M） 2.
记录每个样本在多个 checkpoint 上的 loss → 构建 loss 轨迹向量 3. K-means
聚类（\textasciitilde100 簇）→ 从每簇均匀采样

\textbf{关键结果}： - 11\% MathInstruct 数据匹配全数据性能 -
比之前方法平均高 4.7\%（6 个评测集） - MATH 上 50K 筛选数据提升 Phi-2 达
16.6\%

\textbf{理论基础}：同一 loss 轨迹簇内的样本具有类似梯度 → S2L
子集与全数据集有有界梯度误差。

\textbf{与 SCX 的核心差异}： 1. \textbf{代理模型 vs 状态划分}：S2L
用小模型 loss 轨迹聚类，需要\textbf{先训练一个小模型}。SCX
的状态划分可以用无监督方法（描述符 + 聚类）不需要训练。 2.
\textbf{单轨迹 vs 多专家}：S2L 只用一个代理模型，不利用专家多样性。SCX
天然多专家。 3. \textbf{均匀采样 vs 价值排序}：S2L
在簇内均匀随机采样；SCX 在每个状态内按价值排序后选择。

\textbf{SCX 相对优势}： - S2L 不能处理\textbf{噪声样本问题}------loss
轨迹高的样本既有可能是困难样本也有可能是噪声；SCX 通过噪声分数明确区分。
- S2L 不能利用\textbf{领域先验}（如材料科学中已知的键合状态）；SCX
的状态空间可以融合先验知识。 - S2L 的 K-means 聚类需要预先指定 K
值（\textasciitilde100），SCX 的状态数量可以自适应。

\begin{center}\rule{0.5\linewidth}{0.5pt}\end{center}

\subsubsection{B.4
大规模指令微调数据筛选的其它进展}\label{b.4-ux5927ux89c4ux6a21ux6307ux4ee4ux5faeux8c03ux6570ux636eux7b5bux9009ux7684ux5176ux5b83ux8fdbux5c55}

\paragraph{B.4.1 InstructMining}\label{b.4.1-instructmining}

\begin{itemize}
\tightlist
\item
  \textbf{论文}：Cao et al., 2024
\item
  \textbf{概述}：自动筛选优质指令遵循数据
\end{itemize}

\paragraph{B.4.2 ``Maybe Only 0.5\% Data is
Needed''}\label{b.4.2-maybe-only-0.5-data-is-needed}

\begin{itemize}
\tightlist
\item
  \textbf{论文}：Chen et al., 2023/2024
\item
  \textbf{概述}：极小的高质量子集可能就足够了
\end{itemize}

\paragraph{B.4.3 ``Rethinking Data Selection at Scale: Random Selection
is Almost All You
Need''}\label{b.4.3-rethinking-data-selection-at-scale-random-selection-is-almost-all-you-need}

\begin{itemize}
\tightlist
\item
  \textbf{论文}：2024
\item
  \textbf{核心发现}：大规模下随机选择意外地有竞争力；多样性可能比质量更重要
\end{itemize}

\paragraph{B.4.4 ``A Critical Look at Targeted Instruction Selection''
(Nayak et al.,
2025)}\label{b.4.4-a-critical-look-at-targeted-instruction-selection-nayak-et-al.-2025}

\begin{itemize}
\tightlist
\item
  \textbf{链接}：https://arxiv.org/abs/2602.14696
\item
  \textbf{核心发现}：梯度表示（LESS 风格）最具预测力；贪心 round-robin
  在低预算最优，optimal transport 在高预算最优；随机在足够大规模下仍然强
\end{itemize}

\paragraph{B.4.5 综述论文}\label{b.4.5-ux7efcux8ff0ux8bbaux6587}

\begin{itemize}
\tightlist
\item
  \textbf{``A Survey on Data Selection for LLM Instruction Tuning''}
  (Wang et al., JAIR 2025)
\item
  \textbf{``Unleashing the Power of Data Tsunami''} (Qin et al., TMLR
  2024)
\item
  \textbf{链接}：https://arxiv.org/abs/2408.02085
\end{itemize}

均将方法分为三类：\textbf{基于质量}（困惑度、IFD 分数、GPT
评分）、\textbf{基于多样性}（词汇多样性、k-center
贪心、聚类）、\textbf{基于重要性}（不确定性、loss 梯度、误差）。

\textbf{SCX
补全的视角}：所有现有方法都只考虑数据本身的质量/多样性，没有考虑\textbf{当前模型的状态}和\textbf{专家的条件可靠性}。SCX
的三因素价值公式（误差 + 密度 +
可学习性）是首个同时考虑这三个维度的统一框架。

\begin{center}\rule{0.5\linewidth}{0.5pt}\end{center}

\subsection{C. Active Learning}\label{c.-active-learning}

\subsubsection{C.1 DP-GEN (Deep Potential
Generator)}\label{c.1-dp-gen-deep-potential-generator}

\begin{itemize}
\tightlist
\item
  \textbf{原始论文}：Zhang et al., Phys. Rev.~Materials 3, 023804 (2019)
\item
  \textbf{文档}：https://deepwiki.com/deepmodeling/dpgen/
\item
  \textbf{代码}：https://github.com/deepmodeling/dpgen
\end{itemize}

\textbf{解决问题}：为 Deep Potential
分子动力学模拟自动产生高质量训练数据（DFT 标签）。

\textbf{数学对象}：训练 4 个独立初始化的 DeePMD 模型 ensemble →
计算模型偏差（Model Deviation）作为不确定性：

\begin{verbatim}
σ_f^max = max_i √⟨||f_i - ⟨f_i⟩||²⟩
\end{verbatim}

基于信任阈值（trust level）选择候选结构：σ\_f ∈ {[}lower, upper{]} →
选入；σ\_f \textless{} lower → 丢弃（已准确）；σ\_f \textgreater{} upper
→ 丢弃（太远）。

\textbf{与 SCX 的核心差异}： 1. \textbf{Ensemble vs 状态}：DP-GEN
用多模型 ensemble 的不一致性衡量不确定性；SCX 用状态划分 + 专家可靠性 +
内部一致性。 2. \textbf{统一阈值 vs 状态条件}：DP-GEN
对全局使用统一信任阈值（如 0.12 eV/A）；SCX
的''阈值''是状态条件的------不同状态有不同的接受标准（因为不同状态的内在噪声水平不同）。
3. \textbf{Exploration 驱动 vs 缺陷驱动}：DP-GEN 选 ensemble
不一致的区域（exploration）；SCX 选高误差 - 高密度 -
高一致性的区域（缺陷驱动）。

\textbf{SCX 相对优势}： - DP-GEN 不能区分''模型 ensemble
不一致是因为缺乏数据''还是''因为该区域固有噪声高''。SCX
通过状态内部一致性 C(s) 做区分。 - DP-GEN 需要训练 4
个模型同时推理，计算开销大。SCX 的状态划分只需要一次聚类。 - DP-GEN
的信任阈值是超参数（人工调整）。SCX 的状态价值是自动计算的。

\begin{center}\rule{0.5\linewidth}{0.5pt}\end{center}

\subsubsection{C.2 MTP / MLIP-3 --- Extrapolation Grade +
D-Optimality}\label{c.2-mtp-mlip-3-extrapolation-grade-d-optimality}

\begin{itemize}
\tightlist
\item
  \textbf{工具}：MLIP-2 / MLIP-3 (https://gitlab.com/ashapeev/mlip-2)
\item
  \textbf{方法}：Shapeev, ``Moment Tensor Potentials'', 2016
\end{itemize}

\textbf{解决问题}：为 Moment Tensor Potential (MTP) 的主动学习选择最需要
DFT 标注的新配置。

\textbf{数学对象}： 1. \textbf{Extrapolation Grade γ(cfg)}：对线性
MTP，γ(cfg) = max\_j \textbar c\_j\textbar，其中 c = (b₁\ldots b\_m) ·
A⁻¹，A 是当前 active set 的信息矩阵。若 γ \textgreater{}
1，该配置是外推的（超出当前训练域）。 2. \textbf{D-Optimality
(MaxVol)}：从候选集中选择使信息矩阵行列式最大化的一组配置。

\textbf{工作流}：MD 模拟中若 γ \textgreater{} γ\_th（通常 10），采样所有
γ \textgreater{} 1 的配置 → D-optimality 选子集 → DFT 标注 → 重训练。

\textbf{与 SCX 的核心差异}： 1. \textbf{线性近似 vs 非线性状态}：MTP 的
extrapolation grade 基于线性模型（或线性化的非线性模型）的凸包。SCX
可以处理任意非线性的状态空间。 2. \textbf{几何 vs 语义}：D-optimality
是纯几何准则（最大线性独立），不考虑标签信息。SCX
的状态同时基于特征（描述符）和误差分布。 3. \textbf{单专家 vs
多专家}：MTP 自发主动学习只涉及一个模型。SCX 天然处理多个专家。

\textbf{SCX 相对优势}： - D-optimality
只关注特征空间的几何边界（extrapolation），可能遗漏内插区域的''误差热点''（如相界附近描述符平滑但能量跳变）。SCX
通过 r̄(s) 直接感知误差。 - D-optimality 不能区分噪声和外推区域。SCX
明确区分噪声（N(s)）和可学习困难（L(s)）。 -
近期改进方向（Clustering-enhanced local D-optimality, 2025）试图解决
D-optimality 对稀有机会的遗漏------这正好是 SCX 状态划分擅长的问题。

\begin{center}\rule{0.5\linewidth}{0.5pt}\end{center}

\subsubsection{C.3 DIRECT Sampling (DImensionality-Reduced Encoded
Clusters with
sTratified)}\label{c.3-direct-sampling-dimensionality-reduced-encoded-clusters-with-stratified}

\begin{itemize}
\tightlist
\item
  \textbf{论文}：``DImensionality-Reduced Encoded Clusters with
  sTratified (DIRECT) sampling'', Nature, 2024
\item
  \textbf{链接}：https://www.nature.com/articles/s41524-024-01227-4
\end{itemize}

\textbf{解决问题}：为
MLIP（机器学习势函数）选择代表性训练数据，减少冗余、保持多样性。

\textbf{数学对象}： 1. 将原子环境编码为低维表示 2.
对编码后的表示进行聚类 3. 从每个簇分层采样（stratified sampling）

\textbf{与 SCX 的关系}：DIRECT 的流程（降维 → 聚类 → 分层采样）与 SCX
的''状态划分''思想高度相似。

\textbf{核心差异}： 1. \textbf{纯多样性 vs 误差感知}：DIRECT
只基于描述符空间的距离/密度，不感知模型误差。每个簇中均匀采样，不考虑哪个簇误差更高。SCX
用 r̄(s) 自动识别高误差状态。 2. \textbf{静态 vs 迭代}：DIRECT
是一次性的数据选择。SCX 是迭代主动学习框架------训练 → 评估价值 → 采集 →
再训练。 3. \textbf{无专家 vs 多专家}：DIRECT 不利用专家信息。

\textbf{SCX 优势}：DIRECT 是一个优秀的''多样性基线''方法，可以被集成到
SCX 框架中作为状态划分的候选方法。但 SCX 的误差感知价值函数使它比 DIRECT
更高效------DIRECT 可能为已经准确的区域浪费标注资源。

\begin{center}\rule{0.5\linewidth}{0.5pt}\end{center}

\subsubsection{C.4 DIRECT: Deep Active Learning under Imbalance and
Label
Noise}\label{c.4-direct-deep-active-learning-under-imbalance-and-label-noise}

\begin{itemize}
\tightlist
\item
  \textbf{论文}：Zhang et al., ``DIRECT: Deep Active Learning under
  Imbalance and Label Noise'', 2024
\item
  \textbf{链接}：https://arxiv.org/abs/2312.09196
\end{itemize}

\textbf{解决问题}：在类别不平衡和标签噪声下的深度主动学习。

\textbf{数学对象}：将多类主动学习降维为一系列一维阈值学习问题。通过降维将模型高维输出映射到
1D 分数空间，识别最优分离阈值来选择有信息量的、类别平衡的样本。

\textbf{改进}：比 SOTA 主动学习节省 \textgreater60\%
标注量，支持批量标注，对标签噪声鲁棒。

\textbf{SCX 差异}：这个 DIRECT 是\textbf{分类任务}的主动学习方法，依赖
softmax 输出概率。SCX 可以处理回归任务（如力场拟合）和分类任务。SCX
的噪声处理更一般（不限于标签噪声，还包含状态内不一致性）。

\begin{center}\rule{0.5\linewidth}{0.5pt}\end{center}

\subsubsection{C.5 QActor --- Noise-Aware Active
Sampling}\label{c.5-qactor-noise-aware-active-sampling}

\begin{itemize}
\tightlist
\item
  \textbf{论文}：Ghiassi et al., ``QActor: On-line Active Learning for
  Noisy Labeled Stream Data'', ACML 2021
\item
  \textbf{链接}：https://arxiv.org/abs/2001.10399
\end{itemize}

\textbf{解决问题}：在流式数据场景中，既有噪声标签又有有限标注预算。

\textbf{数学对象}：结合两种策略： 1. 质量模型（Quality
Model）：选择看起来干净的样本 2. 主动查询（Active
Querying）：用标注预算向 oracle 询问最具信息量的样本

在噪声率 30\%-80\% 的数据集上实现接近无噪声训练的精度，最多只需 6\%
额外查询。

\textbf{SCX 差异}： - QActor 的''质量''是全局的（干净 vs
不干净），不区分类别或状态。 - SCX 的状态条件噪声分数（NoiseScore(x\_i)
= r\_i · (1/ρ(s\_i)+ε) · {[}1 -
C(s\_i){]}）可以区分''高误差但可学习''和''高误差且噪声''。 - QActor
面向流式数据，SCX 面向静态池 + 迭代获取。

\begin{center}\rule{0.5\linewidth}{0.5pt}\end{center}

\subsection{D.
竞争格局总结}\label{d.-ux7adeux4e89ux683cux5c40ux603bux7ed3}

\subsubsection{D.1
四象限定位图（文本描述）}\label{d.1-ux56dbux8c61ux9650ux5b9aux4f4dux56feux6587ux672cux63cfux8ff0}

\begin{verbatim}
                  全局/固定                  状态条件
                ┌──────────────────┬──────────────────┐
                │   Data Shapley    │                  │
                │   TMC-Shapley     │     SCX ★       │
    数据估值     │   KNN-Shapley    │     (本工作)      │
                │   Influence Func  │                  │
                │   TRAK / LoGra    │                  │
                ├──────────────────┼──────────────────┤
                │   LESS (梯度相似)  │                  │
                │   S2L (loss轨迹)   │   EGP(误差分区)    │
    数据选择     │   DCLM (过滤)     │   DIRECT(聚类)    │
                │   InstructMining  │   DP-GEN(ensemble)│
                │   Random          │   MTP D-opt      │
                └──────────────────┴──────────────────┘
\end{verbatim}

\textbf{解释}： - \textbf{右上象限}：只有 SCX
同时做到了''状态条件''和''数据价值''的统一。EGP 不涉及数据估值，DIRECT
和 MTP 的 D-optimality 不涉及多专家路由，DP-GEN 不进行状态划分。 -
\textbf{左下象限}：大部分现有工作（Data Shapley、Influence
Functions、LESS、S2L）都在这个象限------它们都是\textbf{全局/固定}的。

\subsubsection{D.2
各方法覆盖的功能矩阵}\label{d.2-ux5404ux65b9ux6cd5ux8986ux76d6ux7684ux529fux80fdux77e9ux9635}

\begin{longtable}[]{@{}lllllll@{}}
\toprule\noalign{}
方法 & 数据估值 & 数据选择 & 多专家 & 状态条件 & 噪声检测 & 可学习性 \\
\midrule\noalign{}
\endhead
\bottomrule\noalign{}
\endlastfoot
Data Shapley & Y & Y & N & N & N & N \\
TMC-Shapley & Y & Y & N & N & N & N \\
KNN-Shapley & Y & Y & N & N & N & N \\
CHG Shapley & Y & Y & N & N & Partial & N \\
Influence Func & Y & Y & N & N & N & N \\
TRAK & Y & Y & N & N & N & N \\
LoGra/LogIX & Y & Y & N & N & N & N \\
LESS & N & Y & N & N & N & N \\
S2L & N & Y & N & N & N & N \\
DCLM & N & Y & N & N & N & N \\
DP-GEN & N & Y & Y(ensemble) & N & N & N \\
MTP D-opt & N & Y & N & N & N & N \\
DIRECT(Nature) & N & Y & N & Y(聚类) & N & N \\
QActor & N & Y & N & N & Y(标签) & N \\
\textbf{SCX} & \textbf{Y} & \textbf{Y} & \textbf{Y} & \textbf{Y} &
\textbf{Y} & \textbf{Y} \\
\end{longtable}

SCX 是唯一在\textbf{所有列}都覆盖的方法。

\subsubsection{D.3
计算复杂度比较}\label{d.3-ux8ba1ux7b97ux590dux6742ux5ea6ux6bd4ux8f83}

\begin{longtable}[]{@{}llll@{}}
\toprule\noalign{}
方法 & 推理复杂度 & 训练复杂度 & 扩展性 \\
\midrule\noalign{}
\endhead
\bottomrule\noalign{}
\endlastfoot
Data Shapley (MC) & O(N² × train) & N/A & 差 (\textless10⁴ 样本) \\
KNN-Shapley & O(n log n) & O(n) & 好 (代理模型) \\
TRAK & O(N × d × k) & O(dk²) & 好 \\
LoGra & O(N × √(dk)) & O(√(dk)) & 很好 \\
LESS & O(N × d × n\_val) & O(d × n\_train) & 中等 \\
S2L & O(N × T) + K-means & O(T × n\_small) & 好 \\
DP-GEN & O(4 × N × fwd) & O(4 × train) & 中等 (4 model ensemble) \\
MTP D-opt & O(N × d²) & O(N × d²) & 好 \\
\textbf{SCX} & \textbf{O(N × d + K × N)} & \textbf{O(N × d + K-means)} &
\textbf{好} \\
\end{longtable}

SCX 的计算主要来自三部分：描述符计算 O(N×d)、状态划分 O(N×d +
K-means)、状态价值计算 O(K×M)。没有模型重训练开销，非常适合大规模问题。

\begin{center}\rule{0.5\linewidth}{0.5pt}\end{center}

\subsection{E. SCX
的独特定位}\label{e.-scx-ux7684ux72ecux7279ux5b9aux4f4d}

\subsubsection{E.1 SCX 的 state-conditioned
视角填补了什么空白？}\label{e.1-scx-ux7684-state-conditioned-ux89c6ux89d2ux586bux8865ux4e86ux4ec0ux4e48ux7a7aux767d}

SCX 填补了机器学习四个基本问题的交叉空白：

\textbf{空白 1：数据价值的条件性}

现有数据估值方法（Data Shapley、Influence
Functions、LOO）都假设数据的''价值''是其固有属性------同一个样本在所有情境下都有相同的价值。SCX
的 Proposition 1 证明这是错的：一个样本在状态 s₁ 可能是关键数据，在状态
s₂ 可能是冗余数据。价值是\textbf{状态条件}的。

这与直觉一致：一个训练好的回归器在密集区域不需要更多数据，但在稀疏或高误差区域需要。现有方法无法表达这个简单直觉。

\textbf{空白 2：专家可靠性的局部性}

现有多专家系统（MoE、集成模型）假设每个专家有一个全局''能力分数''或权重。SCX
证明（Proposition 1 和 3）：专家 A 可能在状态 s₁ 优于专家 B，但在状态 s₂
劣于专家 B。因此专家选择必须状态条件化。

这与材料 MLIP
的实际经验一致：一个在体相材料上训练的势函数可能在表面、缺陷或相界面处完全失效------不是''不好的模型''，而是在那些状态上不适用。

\textbf{空白 3：噪声与困难样本的混淆}

现有主动学习方法（max-residual、uncertainty
sampling）都假设高误差/高不确定性的样本最有价值。SCX 的 Proposition 2
证明：高误差集合 H = H\_learnable ∪
H\_noise。对噪声采样浪费标注预算。必须区分 ``可学习困难的样本'' 和
``固有噪声的样本''。

SCX 通过状态内一致性做区分： - 高误差 + 高密度 + 高一致性 →
可学习困难（宝贵） - 高误差 + 低密度 + 低一致性 → 噪声（应丢弃/降权）

\textbf{空白 4：统一框架}

Data Shapley（估值）、LESS/S2L（LLM
数据选择）、DP-GEN/D-opt（主动学习）、QActor（噪声感知）这四个社区通常各自独立发展。SCX
用一个统一的数学框架（条件风险 + 状态划分 +
专家路由）涵盖了所有四个维度。

\subsubsection{E.2 为什么现有方法无法直接替代
SCX？}\label{e.2-ux4e3aux4ec0ux4e48ux73b0ux6709ux65b9ux6cd5ux65e0ux6cd5ux76f4ux63a5ux66ffux4ee3-scx}

\paragraph{不能替代的理由
1：全局无法表达局部}\label{ux4e0dux80fdux66ffux4ee3ux7684ux7406ux7531-1ux5168ux5c40ux65e0ux6cd5ux8868ux8fbeux5c40ux90e8}

Data Shapley、Influence Functions、LESS、S2L、TRAK
都输出\textbf{全局}分数/选择。如果配置空间存在\textbf{非均匀的误差分布}（MLIP
的核心假设），全局方法必然次优。SCX 的状态划分是解决这个问题的必要机制。

\paragraph{不能替代的理由
2：监督信号未被利用}\label{ux4e0dux80fdux66ffux4ee3ux7684ux7406ux7531-2ux76d1ux7763ux4fe1ux53f7ux672aux88abux5229ux7528}

DCLM、DIRECT、D-optimality
只用到了\textbf{特征空间的结构}（perplexity、距离、密度），没有用到模型误差、标签一致性或专家可靠性。SCX
的估值公式 r̄(s) · ρ(s) · L(s) · {[}1 - D(s){]}
同时利用了特征结构（ρ）和监督信号（r̄, L, SCX\_m）。

\paragraph{不能替代的理由
3：多专家未被整合}\label{ux4e0dux80fdux66ffux4ee3ux7684ux7406ux7531-3ux591aux4e13ux5bb6ux672aux88abux6574ux5408}

Data Shapley、LESS、S2L、DP-GEN 都只有一个模型/代理。DP-GEN 的 ensemble
是\textbf{多个相同模型}（随机种子不同），不是\textbf{不同领域的专家}。SCX
的路由框架可以处理真正异构的专家集合（如 MACE + NEP + ACE
不同类型的势函数）。

\paragraph{不能替代的理由
4：噪声处理不足}\label{ux4e0dux80fdux66ffux4ee3ux7684ux7406ux7531-4ux566aux58f0ux5904ux7406ux4e0dux8db3}

\begin{itemize}
\tightlist
\item
  Data Shapley 和 Influence Functions 没有噪声模型
\item
  LESS 和 S2L 不区分高误差的来源
\item
  DP-GEN 将 ensemble 不一致全部当作''需要数据''而不是''潜在噪声''
\item
  MTP 的 D-optimality 只关注几何外推
\end{itemize}

SCX 是唯一在理论层面明确区分''噪声''和''可学习困难''的方法（Proposition
2 及其证明）。

\paragraph{不能替代的理由
5：没有路由}\label{ux4e0dux80fdux66ffux4ee3ux7684ux7406ux7531-5ux6ca1ux6709ux8defux7531}

SCX 的 expert routing
是其独特组件------不仅选择''要标注什么数据''，还决定''由哪个/哪些专家来标注/预测''。这在
MoE 和联邦学习等场景具有独立价值。

\subsubsection{E.3
与最接近方法的逐一对比}\label{e.3-ux4e0eux6700ux63a5ux8fd1ux65b9ux6cd5ux7684ux9010ux4e00ux5bf9ux6bd4}

\textbf{最接近方法：S2L} -
共同点：都使用聚类划分数据空间，从各类中选代表点 - 差异：S2L
不知道各类的误差有多大，在每类中均匀采样；SCX
知道各类的误差，优先采样高误差类 - 差距：S2L 需要先训练小模型 → 需要 4-7
天；SCX 可以零训练启动（描述符聚类）

\textbf{最接近方法：LESS} - 共同点：都使用''相似性''筛选数据 -
差异：LESS 在梯度空间做相似性（计算昂贵），SCX 在状态空间做划分 -
差距：LESS 不能处理多专家，不能区分噪声

\textbf{最接近方法：DP-GEN} -
共同点：都是迭代主动学习，都基于当前模型的不足选择数据 - 差异：DP-GEN 用
ensemble 不确定性（需要多模型同步），SCX 用状态划分 + 误差估计 -
差距：DP-GEN 不能区分噪声，不能利用领域先验

\textbf{最接近方法：DIRECT (Nature 2024)} - 共同点：都使用降维 +
聚类划分配置空间 - 差异：DIRECT 静态一次选择，不感知误差；SCX 迭代 +
误差感知 - 差距：DIRECT 是全流程相似，但缺了 SCX 的价值函数和迭代机制

\subsubsection{E.4 SCX
的潜在弱点}\label{e.4-scx-ux7684ux6f5cux5728ux5f31ux70b9}

\begin{enumerate}
\def\labelenumi{\arabic{enumi}.}
\tightlist
\item
  \textbf{状态划分质量依赖}：如果状态划分方法（聚类/描述符）不好，后续所有步骤都受影响。需要鲁棒的自动状态发现。
\item
  \textbf{专家先验可用性}：SCX 的多专家优势需要至少有 2
  个可用专家（或一个能分解的专家）。单模型场景下部分功能退化。
\item
  \textbf{小样本状态}：当某个状态内样本数太少时，条件风险 R\_m(s)
  和可学习性 L(s) 的估计不可靠。需要贝叶斯平滑。
\item
  \textbf{状态演化的处理}：随着新数据的加入，状态边界可能变化。需要增量/在线状态更新机制。
\item
  \textbf{计算复杂度}：相比最简单的随机采样或 heuristic 过滤，SCX 的聚类
  + 估值计算开销仍然较高。
\end{enumerate}

\subsubsection{E.5
差异化战略总结}\label{e.5-ux5deeux5f02ux5316ux6218ux7565ux603bux7ed3}

\begin{longtable}[]{@{}lll@{}}
\toprule\noalign{}
SCX 特性 & 不可替代性级别 & 说明 \\
\midrule\noalign{}
\endhead
\bottomrule\noalign{}
\endlastfoot
State-conditioned 价值 & ★★★★★ & 全局方法无法表达状态特定的价值 \\
噪声 vs 困难的区分 & ★★★★★ & 没有其他方法能在理论上做此区分 \\
多专家路由 & ★★★★☆ & 需要有多专家可用场景（材料 MLIP 天然成立） \\
无训练启动 & ★★★★☆ & 只需描述符 + 聚类，比需要训练的方法便宜 \\
统一框架 & ★★★☆☆ & 学术价值高，单一场景不一定需要所有功能 \\
\end{longtable}

\begin{center}\rule{0.5\linewidth}{0.5pt}\end{center}

\subsection{7. 参考文献}\label{ux53c2ux8003ux6587ux732e}

\subsubsection{Data Valuation}\label{data-valuation}

\begin{enumerate}
\def\labelenumi{\arabic{enumi}.}
\tightlist
\item
  Ghorbani \& Zou, ``Data Shapley: Equitable Valuation of Data for
  Machine Learning'', ICML 2019. https://arxiv.org/abs/1904.02868
\item
  Jia et al., ``Towards Efficient Data Valuation Based on the Shapley
  Value'', AISTATS 2019. https://arxiv.org/abs/1902.10275
\item
  Li et al., ``CHG Shapley: Efficient Data Valuation and Selection'',
  2024. https://arxiv.org/abs/2406.11730
\item
  ``Fast-DataShapley: Neural Modeling for Training Data Valuation'',
  2025. https://arxiv.org/abs/2506.05281
\item
  ``Asymmetric Data Shapley for Structure-Aware Valuation'', 2025.
  https://arxiv.org/abs/2511.12863
\item
  ``Geometric Data Valuation via Leverage Scores'', 2025.
  https://arxiv.org/abs/2511.02100
\item
  ``Faithful Group Shapley Value'', NeurIPS 2025.
  https://neurips.cc/virtual/2025/san-diego/poster/115078
\item
  Park et al., ``TRAK: Attributing Model Behavior at Scale'', ICML 2023.
  https://arxiv.org/abs/2212.02522
\item
  Zheng et al., ``D-TRAK: Intriguing Properties of Data Attribution on
  Diffusion Models'', ICLR 2024. https://github.com/qwkkkkk/D-TRAK
\item
  Choe et al., ``What is Your Data Worth to GPT? LLM-Scale Data
  Valuation with Influence Functions'', NeurIPS 2025.
  https://arxiv.org/abs/2405.13954
\item
  ``Generalized Group Data Attribution'', 2024.
  https://arxiv.org/abs/2410.09940
\item
  ``Kairos: Scalable Model-Agnostic Data Valuation'', 2025.
  https://arxiv.org/abs/2506.23799
\end{enumerate}

\subsubsection{LLM Data Selection}\label{llm-data-selection}

\begin{enumerate}
\def\labelenumi{\arabic{enumi}.}
\setcounter{enumi}{12}
\tightlist
\item
  ``DataComp-LM: In search of the next generation of training sets for
  language models'', NeurIPS 2024. https://arxiv.org/abs/2406.11794
\item
  Xia et al., ``LESS: Selecting Influential Data for Targeted
  Instruction Tuning'', ICML 2024. https://arxiv.org/abs/2402.04333
\item
  Yang et al., ``SmallToLarge (S2L): Scalable Data Selection for
  Fine-tuning LLMs'', NeurIPS 2024. https://arxiv.org/abs/2403.07384
\item
  Qin et al., ``Unleashing the Power of Data Tsunami: A Survey on Data
  Assessment and Selection for Instruction Tuning'', TMLR 2024.
  https://arxiv.org/abs/2408.02085
\item
  Nayak et al., ``A Critical Look at Targeted Instruction Selection'',
  2025. https://arxiv.org/abs/2602.14696
\item
  Wang et al., ``A Survey on Data Selection for LLM Instruction
  Tuning'', JAIR 2025.
\end{enumerate}

\subsubsection{Active Learning}\label{active-learning}

\begin{enumerate}
\def\labelenumi{\arabic{enumi}.}
\setcounter{enumi}{18}
\tightlist
\item
  Zhang et al., ``DP-GEN: Deep Potential Generator'', Phys.
  Rev.~Materials 3, 023804 (2019). https://github.com/deepmodeling/dpgen
\item
  Shapeev, ``Moment Tensor Potentials'', 2016.
  https://gitlab.com/ashapeev/mlip-2
\item
  ``DIRECT: DImensionality-Reduced Encoded Clusters with sTratified
  Sampling'', Nature 2024.
  https://www.nature.com/articles/s41524-024-01227-4
\item
  Zhang et al., ``DIRECT: Deep Active Learning under Imbalance and Label
  Noise'', 2024. https://arxiv.org/abs/2312.09196
\item
  Ghiassi et al., ``QActor: On-line Active Learning for Noisy Labeled
  Stream Data'', ACML 2021. https://arxiv.org/abs/2001.10399
\end{enumerate}

\subsubsection{Tools}\label{tools}

\begin{enumerate}
\def\labelenumi{\arabic{enumi}.}
\setcounter{enumi}{23}
\tightlist
\item
  pyDVL (Python Data Valuation Library).
  https://github.com/aai-institute/pyDVL --- v0.10.0, 2024
\item
  LogIX / LoGra. https://github.com/sangkeuncho/LogIX --- 2024
\end{enumerate}

\end{document}
